\begin{theorem}
\label{theorem-vanishing-affine-curves}
If $k$ is an algebraically closed field, $X$ is a separated, finite type
scheme of dimension $\leq 1$ over $k$, and $\mathcal{F}$ is a torsion
abelian sheaf on $X_\etale$, then
\begin{enumerate}
\item
$H^q_\etale(X, \mathcal{F}) = 0$ for $q > 2$,
\item
$H^q_\etale(X, \mathcal{F}) = 0$ for $q > 1$ if $X$ is affine,
\item
$H^q_\etale(X, \mathcal{F}) = 0$ for $q > 1$ if $p = \text{char}(k) > 0$
and $\mathcal{F}$ is $p$-power torsion,
\item
$H^q_\etale(X, \mathcal{F})$ is finite if $\mathcal{F}$ is
constructible and torsion prime to $\text{char}(k)$,
\item
$H^q_\etale(X, \mathcal{F})$ is finite if $X$ is proper and
$\mathcal{F}$ constructible,
\item
$H^q_\etale(X, \mathcal{F}) \to
H^q_\etale(X_{k'}, \mathcal{F}|_{X_{k'}})$ is an isomorphism
for any extension $k'/k$ of algebraically closed fields
if $\mathcal{F}$ is torsion prime to $\text{char}(k)$,
\item
$H^q_\etale(X, \mathcal{F}) \to
H^q_\etale(X_{k'}, \mathcal{F}|_{X_{k'}})$ is an isomorphism
for any extension $k'/k$ of algebraically closed fields
if $X$ is proper,
\item
$H^2_\etale(X, \mathcal{F}) \to H^2_\etale(U, \mathcal{F})$
is surjective for all $U \subset X$ open.
\end{enumerate}
\end{theorem}

\begin{proof}
The theorem says that in Situation \ref{situation-what-to-prove}
statements (\ref{item-vanishing}) -- (\ref{item-surjective}) hold.
Our first step is to replace $X$ by its reduction, which is permissible
by Proposition \ref{proposition-topological-invariance}.
By Lemma \ref{lemma-torsion-colimit-constructible} we can write
$\mathcal{F}$ as a filtered colimit of constructible abelian sheaves.
Taking cohomology commutes with colimits, see Lemma \ref{lemma-colimit}.
Moreover, pullback via $X_{k'} \to X$ commutes with colimits as a left
adjoint. Thus it suffices to prove the statements for a constructible sheaf.

\medskip\noindent
In this paragraph we use Lemma \ref{lemma-ses-statements} without further
mention. Writing
$\mathcal{F} = \mathcal{F}_1 \oplus \ldots \oplus \mathcal{F}_r$
where $\mathcal{F}_i$ is $\ell_i$-primary for some prime $\ell_i$, we may
assume that $\ell^n$ kills $\mathcal{F}$ for some prime $\ell$. Now consider
the exact sequence
$$
0 \to \mathcal{F}[\ell] \to \mathcal{F} \to \mathcal{F}/\mathcal{F}[\ell] \to 0.
$$
Thus we see that it suffices to assume that $\mathcal{F}$ is $\ell$-torsion.
This means that $\mathcal{F}$ is a constructible sheaf of
$\mathbf{F}_\ell$-vector spaces for some prime number $\ell$.

\medskip\noindent
By definition this means there is a dense open $U \subset X$
such that $\mathcal{F}|_U$ is finite locally constant sheaf of
$\mathbf{F}_\ell$-vector spaces. Since $\dim(X) \leq 1$ we may
assume, after shrinking $U$, that $U = U_1 \amalg \ldots \amalg U_n$
is a disjoint union of irreducible schemes (just remove the closed
points which lie in the intersections of $\geq 2$ components of $U$).
By Lemma \ref{lemma-restrict-to-open} we reduce to the case
$\mathcal{F} = j_!\mathcal{G}$ where $\mathcal{G}$ is a finite
locally constant sheaf of $\mathbf{F}_\ell$-vector spaces on $U$.

\medskip\noindent
Since we chose $U = U_1 \amalg \ldots \amalg U_n$ with $U_i$ irreducible
we have
$$
j_!\mathcal{G} =
j_{1!}(\mathcal{G}|_{U_1}) \oplus \ldots \oplus
j_{n!}(\mathcal{G}|_{U_n})
$$
where $j_i : U_i \to X$ is the inclusion morphism.
The case of $j_{i!}(\mathcal{G}|_{U_i})$ is handled in
Lemma \ref{lemma-vanishing-easier}.
\end{proof}